
on
définit un @@M0@@-homomorphisme : @@M1@@ prolongeant l'augmentation de @@M2@@ en envoyant @@M3@@ sur @@M4@@, ce qui munit @@M5@@ d'une structure de @@M6@@-@@M7@@-algèbre augmentée. On se propose de déterminer sa filtration (3.10) en fonction de celle de @@M8@@.

On gradue @@M9@@ par le poids par rapport aux @@M10@@, @@M11@@ étant affecté du poids @@M12@@~; on désigne par @@M13@@ l'ensemble des monômes de poids @@M14@@ par rapport aux indéterminées @@M15@@.

\textbf{Lemme 4.11.} \uline{Soit} @@M16@@~; \uline{le coefficient d'un monôme de poids} @@M17@@ \uline{par rapport} \uline{aux} @@M18@@ \uline{dans} @@M19@@ \uline{est de filtration au moins} @@M20@@ (\uline{dans} @@M21@@)

Le polynôme @@M22@@ est sans composantes isobares de poids @@M23@@ par rapport aux @@M24@@, @@M25@@ (3.5) donc le lemme est vrai si @@M26@@ est d'augmentation nulle.

Si @@M27@@ désigne l'homomorphisme d'augmentation : @@M28@@, on peut écrire dans le cas général :
@@M29@@
Or @@M30@@, les @@M31@@ étant isobares de poids @@M32@@ par rapport aux @@M33@@ (3.8)~; on en déduit le lemme en faisant le produit.

\textbf{Proposition 4.12.} @@M34@@.

Il est clair que le second idéal est contenu dans le premier~; il y a donc simplement à montrer la réciproque.

Comme la filtration définie par la famille des idéaux du second membre pour @@M35@@ variable est une filtration d'anneau, il suffit de montrer la réciproque pour un élément du type @@M36@@, avec @@M37@@, ces éléments
